% abstract
\thispagestyle{empty}
\section*{Abstract}
In recent years various authors have developed a new numerical approach to cosmological simulations that formulates the equations describing large scale structure (LSS) formation within a quantum mechanical framework. This method couples the Schr\"odinger and Poisson equations. Previously, work has evolved mainly along two different strands of thought: (1) solving the full system of equations as Widrow \& Kaiser attempted, (2) as an approximation to the full set of equations (the Free Particle Approximation developed by Coles, Spencer and Short). It has been suggested that this approach can be considered in two ways: (1) as a purely classical system that includes more physics than just gravity, or (2) as the representation of a dark matter field, perhaps an Axion field, where the de Broglie wavelength of the particles is large.

In the quasi-linear regime, the Free Particle Approximation (FPA) is amenable to exact solution via standard techniques from the quantum mechanics literature. However, this method breaks down in the fully non-linear regime when shell crossing occurs (confer the Zel'dovich approximation). The first eighteen months of my PhD involved investigating the performance of illustrative 1-D and 3-D ``toy" models, as well as a test against the 3-D code Hydra. Much of this work is a reproduction of the work of Short, and I was able to verify and confirm his results. As an extension to his work I introduced a way of calculating the velocity via the probability current rather than using a phase unwrapping technique. Using the probability current deals directly with the wavefunction and provides a faster method of calculation in three dimensions.

After working on the FPA I went on to develop a cosmological code that did not approximate the Schr\"odinger-Poisson system. The final code considered the full Schr\"odinger equation with the inclusion of a self-consistent gravitational potential via the Poisson equation. This method follows on from Widrow \& Kaiser but extends their method from 2D to 3D, it includes periodic boundary conditions, and cosmological expansion. Widrow \& Kaiser provided expansion via a change of variables in their Schr\"odinger equation; however, this was specific only to the Einstein-de Sitter model. In this thesis I provide a generalization of that approach which works for any flat universe that obeys the Robertson-Walker metric.

In this thesis I aim to provide a comprehensive review of the FPA and of the Widrow-Kaiser method. I hope this work serves as an easy first point of contact to the wave-mechanical approach to LSS and that this work also serves as a solid reference point for all future research in this new field.
